\section{Noise-aware injection scheme}
\label{sec:improved-injection-scheme}

In this section, we modify the scheme in \Cref{sec:k-injection} to improve the injection error rate via post-selection, inspired by Ref.~\cite{li_magic_2015}. The key idea is to directly prepare the state in the $+1$ eigenstate of the desired logical observables $ \bar M_i $,
under the \emph{original} logical representatives $\mathcal{L}$.
This method has two advantages: it removes the use of Clifford rotation, and enables post-selection using stabilizers with deterministic outcome in the first SE round. We call these \emph{fixed stabilizers}. Once an error flips the measurement outcome of any fixed stabilizer, we discard the shot and thus decrease the logical error rate of the kept shots.
Such fixed stabilizers do not exist in the general $k$-injection scheme since all the stabilizers in $\mathcal{S}$ anticommute with the prepared single-qubit Pauli states by construction.

An example of injecting two logicals $i$ and $j$ using the improved injection scheme is shown in \Cref{fig:inject-two-logicals}(d).
In the following, we first introduce sufficient conditions to prepare the target state.
Then, we introduce how to post-select on the measurement outcomes of fixed stabilizers to improve the injection error rate.

\begin{figure*}[htb]
    \centering
    \includegraphics[width=\textwidth]{figs/inject_two_logicals.pdf}
    \caption{High-level summary of this work.
    (a) In-logical overlap: Typical injection scheme for a single logical qubit.
    For surface code, we perform $ R_z(\theta) $ on the injection site $ Q_i $,
    which is the single-qubit sitting in the overlap of its logical $ X $ and $ Z $ supports. But for qLDPC codes, $ O_{i} := \supp(\bar X_i) \cap \supp(\bar Z_i) $ has cardinality above 1.
    The solution is to prepare Bell states on $ O_{i} \setminus \{Q_i\} $.
    (b) Cross-logical overlap (remainder overlap): When injecting two logical qubits, $ i $ and $ j $, the overlap of their supports, i.e., $ C_{i,j} $, may not be empty.
    The solution is to prepare Bell states on $ C_{i,j} $.
    (c) Overlap of overlaps: the above-mentioned types of overlaps will overlap with each other.
    The solution is to identify a \emph{compatible pairing} (\Cref{def:generic-compatible}) s.t. the overlaps can be filled with Bell states.
    (d) An example of injecting two logical qubits, $ i $ and $ j $.
    Prepare Bell states on the overlaps,
    and initialize the rest of the qubits in the appropriate Pauli eigenbases.
    % \Shifan{I would suggest changing X support to something like physical qubits supporting X logical operator, to make it more intuitive.}
    % \Kun{Tried but visually tedious. I think journal readers would be more familiar with the term ``support''?}
    }
    \label{fig:inject-two-logicals}
\end{figure*}


The challenge in preparing the desired eigenstate of $ \bar M_i $'s is to deal with the overlaps of the logical $ X $ and $ Z $ supports due to the nature of qLDPC codes.
The central idea is to resolve these overlaps by Bell state preparation on carefully chosen qubit pairs.

We formalize as follows three overlap types, in-logical overlap, cross-logical overlap, and overlap of overlaps, as shown in \Cref{fig:inject-two-logicals}(a,b,c).
For each injected logical index $i\in\cI\subseteq [k]$, define the in-logical overlap $O_i$
as the overlap of the logical $ X $ and $ Z $ supports:
\begin{equation}
O_i \;:=\; \supp(\bar X_i)\cap\supp(\bar Z_i).
\end{equation}
The \emph{$X$-remainder} and \emph{$Z$-remainder} are defined as the qubits outside the in-logical overlap:
\begin{equation}
A_i \;:=\; \supp(\bar X_i)\setminus O_i, \quad
B_i \;:=\; \supp(\bar Z_i)\setminus O_i.
\end{equation}
Because $\{\bar X_i,\bar Z_i\}=0$ in CSS codes, $|O_i|$ is odd.
Choose an \emph{injection site} $Q_i\in O_i$, where the qubit $Q_i$ is prepared in the $+1$ eigenstate of $M_{Q_i}$. The remaining qubits in the in-logical overlap are denoted as $O_i\setminus\{Q_i\}$.
To treat in-logical and cross-logical overlaps uniformly, we define
\begin{equation}
    C_{i,j} =
    \begin{cases}
        A_i\cap B_j, & i\neq j,\\
        O_i\setminus\{Q_i\}, & i=j.
    \end{cases}
    \label{eq:unified-Cij}
\end{equation}
and
\begin{equation}
C := \bigcup_{i,j\in\cI} C_{i,j}.
\end{equation}

The goal of preparing
the state in the $+1$ eigenstate of the desired logical rotated Pauli observables $ \set{\bar M_i} $ is to satisfy the following logical constraints:
\begin{equation}
\bar M_i|\psi_{\mathrm{in}}\rangle = |\psi_{\mathrm{in}}\rangle,\quad \forall i\in\cI.
\label{eq:goal-pre-se-state}
\end{equation}

Since $[\bar M_i,\proj_{\mathcal S}]=0$, projection by ideal SE preserves the logical constraints:
\begin{equation}
\bar M_i|\psi_{\mathrm{out}}\rangle = |\psi_{\mathrm{out}}\rangle,\quad \forall i\in\cI.
\end{equation}

The goal in \Cref{eq:goal-pre-se-state} holds
if the following two \emph{sufficient} conditions hold:

\begin{definition}[Injection site independence]
    \label{def:injection-site-independence}
\begin{equation}
Q_i \notin \supp(\bar X_j)\cup \supp(\bar Z_j) \qquad \forall i\neq j \text{ in }\cI.
\end{equation}
\end{definition}

\begin{definition}[Compatible pairing]
    \label{def:generic-compatible}
    A pairing $\cP=\{\{u_t,v_t\}\}_{t=1}^T$ of $C$ is \emph{compatible} if for every $i,j\in\cI$ and every pair $\{u,v\}\in\cP$,
    \begin{equation}
    u\in C_{i,j}\ \Longleftrightarrow\ v\in C_{i,j}.
    \end{equation}
    Equivalently, each pair is either fully inside or disjoint from every $C_{i,j}$.
\end{definition}

\begin{theorem}[Sufficient conditions for improved injection scheme]
\label{thm:unified-sufficient-condition}
Assume injection site independence (\Cref{def:injection-site-independence}) and a compatible pairing (\Cref{def:generic-compatible}).
Prepare the state by:
\begin{enumerate}
    \item Bell state on every pair in $\cP$;
    \item $\ket{+}$ on $A_i\setminus C$ for each $i\in\cI$;
    \item $\ket{0}$ on $B_i\setminus C$ for each $i\in\cI$;
    \item the $+1$ eigenstate of $M_{Q_i}$ on each $Q_i$.
\end{enumerate}
Then \Cref{eq:goal-pre-se-state} holds, i.e.,
$
\bar M_i|\psi_{\mathrm{in}}\rangle = |\psi_{\mathrm{in}}\rangle
$, $\forall i\in\cI$.
\end{theorem}
The proof is given in \Cref{sec:proof-all-in-one}. To satisfy injection site independence and compatible pairing,
we might not be able to inject the full set of $ k $ logicals.
We discuss how to find maximal injectable sets in \Cref{sec:maximal-injectable-sets}.

% \subsection{Reducing injection error via fixed-stabilizer post-selection}
% \label{subsec:improve-injection-error-rate}

After fixing the preparation on qubits in
$\supp(\bar X_i)\cup \supp(\bar Z_i)$ for all injected logicals,
the remaining design freedom lies in how we initialize qubits outside these supports.
We use this freedom to increase the protection of the injected logical observables in the first syndrome extraction round.

The key mechanism is post-selection on \emph{fixed stabilizers}.
Our initialization pattern ensures that in the absence of any noise, all the fixed stabilizers deterministically yield a $+1$ outcome.
In the presence of noise, if any fixed stabilizer is measured as $-1$, we discard the shot.
Thus, the usefulness of a given initialization pattern is determined by which first-order error mechanisms can be detected through flipped fixed stabilizers.

Our objective is to choose the remaining qubit initializations so as to maximize the fixed-stabilizer protection of the injected logical supports.
To keep the search space manageable and to avoid introducing additional first-round complexities, we restrict these extra choices to single-qubit Pauli-basis preparations.

For qLDPC codes, this optimization problem is substantially more challenging than in surface code injection schemes.
% qLDPC stabilizers are typically higher weight, lack simple 2D geometric structure, and act on a larger number of physical qubits.
In \Cref{sec:max-protection-injected-supports}, we formulate the search for a high-protection initialization pattern as a mixed-integer linear programming (MILP) problem and solve it with off-the-shelf solvers.
